% Source PDF: source/普通高考/2015/2015山东文.pdf, page 1.
% PDF embedded-bitmap attempt: pdfimages found no image objects; the chart is vector line art.
% Grid evidence: tmp/review_flowchart_2015_shandong_liberal/source300/page1_grid100.jpg
% and q11_original_grid50.jpg; comparison source is q11_original.png from the no-grid page render.
% Elements: rounded start/end terminals, input/output parallelograms, decision diamond,
% process boxes, directed connectors, true/false labels, and the left loop-back branch.
% Directed edges: start->input->test; test-(是)->xplus->test (left loop);
% test-(否)->ysquare->output->finish. No xplus->output edge and no dashed segments.
\begin{tikzpicture}[
  >=Stealth,
  line width=0.45pt,
  line cap=round,
  line join=round,
  font=\normalsize,
  flowbox/.style={draw,anchor=center,align=center,
    execute at begin node={\setlength{\parindent}{0pt}},
    minimum height=0.68cm,inner xsep=4pt,inner ysep=2pt},
  terminal/.style={flowbox,rounded corners=2.5pt,minimum width=1.25cm,minimum height=0.78cm},
  process/.style={flowbox},
  io/.style={flowbox,trapezium,trapezium left angle=72,trapezium right angle=108},
  branchlabel/.style={anchor=center,align=center,inner sep=0pt,
    execute at begin node={\setlength{\parindent}{0pt}}},
  arrow/.style={-{Stealth[length=2.2mm,width=1.6mm]},line width=0.45pt}
]
  % Center column and two branches follow the source PDF's schematic topology.
  \node[terminal] (start) at (0,0) {开始};
  \node[io,minimum width=1.90cm,minimum height=0.82cm] (input) at (0,-1.10) {输入 $x$};
  \node[draw,diamond,aspect=2.7,minimum width=3.25cm,minimum height=1.05cm,
    inner sep=0pt] (test) at (0,-2.38) {$x<2$};
  \node[process,minimum width=2.60cm,minimum height=0.85cm] (xplus) at (0,-3.58) {$x=x+1$};
  \node[process,minimum width=3.40cm,minimum height=1.05cm] (ysquare) at (3.40,-3.58) {$y=3x^2+1$};
  \node[io,minimum width=2.00cm,minimum height=0.82cm] (output) at (3.40,-4.78) {输出 $y$};
  \node[terminal] (finish) at (3.40,-5.92) {结束};

  % Main/true path; loop-back returns to the single decision entry.
  \draw[arrow] (start.south) -- (input.north);
  \draw[arrow] (input.south) -- (test.north);
  \draw[arrow] (test.south) -- (xplus.north);
  \draw (xplus.south) -- (0,-4.15) -- (-2.95,-4.15) -- (-2.95,-1.70);
  % The horizontal loop arrow merges into the existing vertical decision stem;
  % only the main stem carries the downward entry arrowhead.
  \draw[arrow] (-2.95,-1.70) -- (0,-1.70);

  % False path: right branch enters y=3x^2+1, then output and finish.
  % The downward arrow terminates exactly at the process box's north border.
  \draw[arrow] (test.east) -- (3.40,-2.38) -- (ysquare.north);
  \draw[arrow] (ysquare.south) -- (output.north);
  \draw[arrow] (output.south) -- (finish.north);

  \node[branchlabel] at (0.45,-3.00) {是};
  \node[branchlabel] at (1.35,-2.02) {否};
\end{tikzpicture}
